% --- Archon physics formalization source begin ---
% archon:physics
% archon:covers PhyXMiniProblems/problem_phyx_mini_0492.lean
% archon:source-report reports/phyx_mini/problem_phyx_mini_0492.source.json

\chapter{Physics problem phyx\_mini\_0492}
\label{ch:PhyXMiniProblems_problem_phyx_mini_0492}

\paragraph{Problem source.}
\#\# Physical scenario

Consider the following two-step process. Heat is allowed to flow out of an ideal gas at constant volume so that its pressure drops from \textbackslash{}(2.2\textbackslash{},\textbackslash{}mathrm\{atm\}\textbackslash{}) to \textbackslash{}(1.4\textbackslash{},\textbackslash{}mathrm\{atm\}\textbackslash{}). Then the gas expands at constant pressure, from a volume of \textbackslash{}(5.9\textbackslash{},\textbackslash{}mathrm\{L\}\textbackslash{}) to \textbackslash{}(9.3\textbackslash{},\textbackslash{}mathrm\{L\}\textbackslash{}), where the temperature reaches its original value.

\#\# Question

Calculate the total work done by the gas in the process.

\#\# Answer choices

- A: A: 390J

- B: B: 420J

- C: C: 450J

- D: D: 480J

\#\# Figure caption (auxiliary; use the image as primary evidence)

The image is a graph depicting a pressure-volume (P-V) relationship. Here's a detailed description:

- **Axes**: 
  - The y-axis represents pressure (P) in atm, labeled from 0 to above 2 atm.
  - The x-axis represents volume (V) in liters (L), labeled from 0 to beyond 9 L.

- **Curve**: 
  - A red curve shows the relationship between P and V, decreasing as volume increases.

- **Points and Lines**:
  - Point **a** is on the curve at 2.2 atm with a vertical dashed line extending to the curve.
  - Point **b** is at 1.4 atm, with horizontal and vertical dashed lines intersecting the curve at around 5.9 L.
  - Point **c** is on the curve at 1.4 atm, intersecting with a vertical dashed line at 9.3 L.

- **Arrows**:
  - A red arrow indicates a vertical transition from point a to point b.
  - A red arrow indicates a horizontal transition from point b to point c.

- **Annotations**:
  - Points a, b, and c are labeled next to the respective points on the graph.

\#\# Dataset metadata

Laws of Thermodynamics; Physical Model Grounding Reasoning, Multi-Formula Reasoning

\paragraph{Recorded answer/context.}
D: D: 480J

\paragraph{Figure/image path.}
/root/proposal\_for\_physic/science-mango/phyx\_mini\_run/phyx\_data/test\_image/492.png

\paragraph{Formalization target.}
create a compiling Lean file with sorry bodies at `PhyXMiniProblems/problem\_phyx\_mini\_0492.lean`. The Lean declarations must preserve the physical quantities, dimensions or dimensional roles, figure labels, governing-law hypotheses, and final relation expressed by this problem.
Use Mathlib/Physlib names found through LeanExplore where available. If a physics API is missing, introduce faithful local abstractions rather than scalar placeholder aliases.

\begin{theorem}[Physics formalization target]
\label{thm:physics:phyx_mini_0492:target}
\lean{PhyXMiniProblems.ProblemPhyXMini0492.totalWorkDoneByGas_matches_recordedAnswerD}
\uses{def:physics:phyx-mini-0492:phyxminiproblems-problemphyxmini0492-energyinjoules, def:physics:phyx-mini-0492:phyxminiproblems-problemphyxmini0492-twostepidealgasprocess, def:physics:phyx-mini-0492:phyxminiproblems-problemphyxmini0492-matchesproblemstatement, def:physics:phyx-mini-0492:phyxminiproblems-problemphyxmini0492-matchesprimarypressurevolumefigure, def:physics:phyx-mini-0492:phyxminiproblems-problemphyxmini0492-hasphysicalthermodynamicparameters, def:physics:phyx-mini-0492:phyxminiproblems-problemphyxmini0492-satisfiesboundaryworklaws, def:physics:phyx-mini-0492:phyxminiproblems-problemphyxmini0492-answerchoice, def:physics:phyx-mini-0492:phyxminiproblems-problemphyxmini0492-recordedanswerchoice, def:physics:phyx-mini-0492:phyxminiproblems-problemphyxmini0492-roundstonearesttenjoules}
The assigned autoformalize agent should translate this physics problem into Lean declarations in the covered file, with theorem and lemma proof bodies written as `by sorry`.
\end{theorem}
\begin{proof}
This is an autoformalization task, not a proof task. Produce statements that can later be proved by the physics prover without weakening the physical model.
\end{proof}

% --- Archon declaration topology begin ---
\section{Declaration topology}
\begin{definition}[Volume Quantity]
  \label{def:physics:phyx-mini-0492:phyxminiproblems-problemphyxmini0492-volumequantity}
  \lean{PhyXMiniProblems.ProblemPhyXMini0492.VolumeQuantity}
  A nonnegative physical gas volume carrying dimension L³
\end{definition}
\begin{proof}
  This is the declared model definition of Volume Quantity.
\end{proof}
\begin{definition}[pressure In Pascals]
  \label{def:physics:phyx-mini-0492:phyxminiproblems-problemphyxmini0492-pressureinpascals}
  \lean{PhyXMiniProblems.ProblemPhyXMini0492.pressureInPascals}
  Read a physical pressure in coherent SI pascals
\end{definition}
\begin{proof}
  This is the declared model definition of pressure In Pascals.
\end{proof}
\begin{definition}[pressure In Atmospheres]
  \label{def:physics:phyx-mini-0492:phyxminiproblems-problemphyxmini0492-pressureinatmospheres}
  \lean{PhyXMiniProblems.ProblemPhyXMini0492.pressureInAtmospheres}
  \uses{def:physics:phyx-mini-0492:phyxminiproblems-problemphyxmini0492-pressureinpascals}
  Read a pressure in the standard atmospheres printed on the diagram
\end{definition}
\begin{proof}
  This is the declared model definition of pressure In Atmospheres.
\end{proof}
\begin{definition}[volume In Cubic Meters]
  \label{def:physics:phyx-mini-0492:phyxminiproblems-problemphyxmini0492-volumeincubicmeters}
  \lean{PhyXMiniProblems.ProblemPhyXMini0492.volumeInCubicMeters}
  \uses{def:physics:phyx-mini-0492:phyxminiproblems-problemphyxmini0492-volumequantity}
  Read a physical volume in coherent SI cubic metres
\end{definition}
\begin{proof}
  This is the declared model definition of volume In Cubic Meters.
\end{proof}
\begin{definition}[volume In Liters]
  \label{def:physics:phyx-mini-0492:phyxminiproblems-problemphyxmini0492-volumeinliters}
  \lean{PhyXMiniProblems.ProblemPhyXMini0492.volumeInLiters}
  \uses{def:physics:phyx-mini-0492:phyxminiproblems-problemphyxmini0492-volumequantity, def:physics:phyx-mini-0492:phyxminiproblems-problemphyxmini0492-volumeincubicmeters}
  Read a physical volume in the litres printed on the diagram
\end{definition}
\begin{proof}
  This is the declared model definition of volume In Liters.
\end{proof}
\begin{definition}[energy In Joules]
  \label{def:physics:phyx-mini-0492:phyxminiproblems-problemphyxmini0492-energyinjoules}
  \lean{PhyXMiniProblems.ProblemPhyXMini0492.energyInJoules}
  Read signed heat or work in coherent SI joules
\end{definition}
\begin{proof}
  This is the declared model definition of energy In Joules.
\end{proof}
\begin{definition}[joules Per Liter Atmosphere]
  \label{def:physics:phyx-mini-0492:phyxminiproblems-problemphyxmini0492-joulesperliteratmosphere}
  \lean{PhyXMiniProblems.ProblemPhyXMini0492.joulesPerLiterAtmosphere}
  \uses{def:physics:phyx-mini-0492:phyxminiproblems-problemphyxmini0492-pressureinpascals}
  The exact number of joules represented by one litre-atmosphere
\end{definition}
\begin{proof}
  This is the declared model definition of joules Per Liter Atmosphere.
\end{proof}
\begin{definition}[Gas Model]
  \label{def:physics:phyx-mini-0492:phyxminiproblems-problemphyxmini0492-gasmodel}
  \lean{PhyXMiniProblems.ProblemPhyXMini0492.GasModel}
  Equation-of-state model specified by the phrase “ideal gas”
\end{definition}
\begin{proof}
  This is the declared model definition of Gas Model.
\end{proof}
\begin{definition}[Figure Point]
  \label{def:physics:phyx-mini-0492:phyxminiproblems-problemphyxmini0492-figurepoint}
  \lean{PhyXMiniProblems.ProblemPhyXMini0492.FigurePoint}
  The three point labels printed in the pressure--volume diagram
\end{definition}
\begin{proof}
  This is the declared model definition of Figure Point.
\end{proof}
\begin{definition}[Process Leg]
  \label{def:physics:phyx-mini-0492:phyxminiproblems-problemphyxmini0492-processleg}
  \lean{PhyXMiniProblems.ProblemPhyXMini0492.ProcessLeg}
  The two directed process legs shown by red arrows
\end{definition}
\begin{proof}
  This is the declared model definition of Process Leg.
\end{proof}
\begin{definition}[initial Point]
  \label{def:physics:phyx-mini-0492:phyxminiproblems-problemphyxmini0492-processleg-initialpoint}
  \lean{PhyXMiniProblems.ProblemPhyXMini0492.ProcessLeg.initialPoint}
  \uses{def:physics:phyx-mini-0492:phyxminiproblems-problemphyxmini0492-figurepoint, def:physics:phyx-mini-0492:phyxminiproblems-problemphyxmini0492-processleg}
  Initial endpoint of each directed process leg
\end{definition}
\begin{proof}
  This is the declared model definition of initial Point.
\end{proof}
\begin{definition}[final Point]
  \label{def:physics:phyx-mini-0492:phyxminiproblems-problemphyxmini0492-processleg-finalpoint}
  \lean{PhyXMiniProblems.ProblemPhyXMini0492.ProcessLeg.finalPoint}
  \uses{def:physics:phyx-mini-0492:phyxminiproblems-problemphyxmini0492-figurepoint, def:physics:phyx-mini-0492:phyxminiproblems-problemphyxmini0492-processleg}
  Final endpoint of each directed process leg
\end{definition}
\begin{proof}
  This is the declared model definition of final Point.
\end{proof}
\begin{definition}[Process Kind]
  \label{def:physics:phyx-mini-0492:phyxminiproblems-problemphyxmini0492-processkind}
  \lean{PhyXMiniProblems.ProblemPhyXMini0492.ProcessKind}
  Thermodynamic constraint imposed on each process leg
\end{definition}
\begin{proof}
  This is the declared model definition of Process Kind.
\end{proof}
\begin{definition}[Segment Geometry]
  \label{def:physics:phyx-mini-0492:phyxminiproblems-problemphyxmini0492-segmentgeometry}
  \lean{PhyXMiniProblems.ProblemPhyXMini0492.SegmentGeometry}
  Geometric appearance of a process leg in the pressure--volume plane
\end{definition}
\begin{proof}
  This is the declared model definition of Segment Geometry.
\end{proof}
\begin{definition}[Axis Quantity]
  \label{def:physics:phyx-mini-0492:phyxminiproblems-problemphyxmini0492-axisquantity}
  \lean{PhyXMiniProblems.ProblemPhyXMini0492.AxisQuantity}
  Physical quantity assigned to one of the plot axes
\end{definition}
\begin{proof}
  This is the declared model definition of Axis Quantity.
\end{proof}
\begin{definition}[Axis Unit]
  \label{def:physics:phyx-mini-0492:phyxminiproblems-problemphyxmini0492-axisunit}
  \lean{PhyXMiniProblems.ProblemPhyXMini0492.AxisUnit}
  Unit printed beside an axis of the diagram
\end{definition}
\begin{proof}
  This is the declared model definition of Axis Unit.
\end{proof}
\begin{definition}[Comparison Curve Kind]
  \label{def:physics:phyx-mini-0492:phyxminiproblems-problemphyxmini0492-comparisoncurvekind}
  \lean{PhyXMiniProblems.ProblemPhyXMini0492.ComparisonCurveKind}
  The red comparison curve through a and c
\end{definition}
\begin{proof}
  This is the declared model definition of Comparison Curve Kind.
\end{proof}
\begin{definition}[Energy Sign Convention]
  \label{def:physics:phyx-mini-0492:phyxminiproblems-problemphyxmini0492-energysignconvention}
  \lean{PhyXMiniProblems.ProblemPhyXMini0492.EnergySignConvention}
  Sign conventions used for the directed energy transfers
\end{definition}
\begin{proof}
  This is the declared model definition of Energy Sign Convention.
\end{proof}
\begin{definition}[Thermodynamic State]
  \label{def:physics:phyx-mini-0492:phyxminiproblems-problemphyxmini0492-thermodynamicstate}
  \lean{PhyXMiniProblems.ProblemPhyXMini0492.ThermodynamicState}
  \uses{def:physics:phyx-mini-0492:phyxminiproblems-problemphyxmini0492-volumequantity}
  Pressure, volume, and absolute temperature at one labeled state
\end{definition}
\begin{proof}
  This is the declared model definition of Thermodynamic State.
\end{proof}
\begin{definition}[Pressure Volume Figure]
  \label{def:physics:phyx-mini-0492:phyxminiproblems-problemphyxmini0492-pressurevolumefigure}
  \lean{PhyXMiniProblems.ProblemPhyXMini0492.PressureVolumeFigure}
  \uses{def:physics:phyx-mini-0492:phyxminiproblems-problemphyxmini0492-figurepoint, def:physics:phyx-mini-0492:phyxminiproblems-problemphyxmini0492-processleg, def:physics:phyx-mini-0492:phyxminiproblems-problemphyxmini0492-segmentgeometry, def:physics:phyx-mini-0492:phyxminiproblems-problemphyxmini0492-axisquantity, def:physics:phyx-mini-0492:phyxminiproblems-problemphyxmini0492-axisunit, def:physics:phyx-mini-0492:phyxminiproblems-problemphyxmini0492-comparisoncurvekind}
  Qualitative content transcribed from the primary bitmap. Numerical state coordinates remain in the physical states and are recorded separately by MatchesPrimaryPressureVolumeFigure
\end{definition}
\begin{proof}
  This is the declared model definition of Pressure Volume Figure.
\end{proof}
\begin{definition}[Two Step Ideal Gas Process]
  \label{def:physics:phyx-mini-0492:phyxminiproblems-problemphyxmini0492-twostepidealgasprocess}
  \lean{PhyXMiniProblems.ProblemPhyXMini0492.TwoStepIdealGasProcess}
  \uses{def:physics:phyx-mini-0492:phyxminiproblems-problemphyxmini0492-gasmodel, def:physics:phyx-mini-0492:phyxminiproblems-problemphyxmini0492-figurepoint, def:physics:phyx-mini-0492:phyxminiproblems-problemphyxmini0492-processleg, def:physics:phyx-mini-0492:phyxminiproblems-problemphyxmini0492-processkind, def:physics:phyx-mini-0492:phyxminiproblems-problemphyxmini0492-energysignconvention, def:physics:phyx-mini-0492:phyxminiproblems-problemphyxmini0492-thermodynamicstate, def:physics:phyx-mini-0492:phyxminiproblems-problemphyxmini0492-pressurevolumefigure}
  Independent physical quantities for the same closed gas sample. In particular, totalWorkDoneByGas is an independent dimensionful observable; it is not defined to be the requested numerical answer
\end{definition}
\begin{proof}
  This is the declared model definition of Two Step Ideal Gas Process.
\end{proof}
\begin{definition}[Matches Problem Statement]
  \label{def:physics:phyx-mini-0492:phyxminiproblems-problemphyxmini0492-matchesproblemstatement}
  \lean{PhyXMiniProblems.ProblemPhyXMini0492.MatchesProblemStatement}
  \uses{def:physics:phyx-mini-0492:phyxminiproblems-problemphyxmini0492-energyinjoules, def:physics:phyx-mini-0492:phyxminiproblems-problemphyxmini0492-twostepidealgasprocess}
  Facts stated in the prose. The strict negative heat readout encodes heat flow out of the gas under the declared sign convention. No work value or answer choice occurs here
\end{definition}
\begin{proof}
  This is the declared model definition of Matches Problem Statement.
\end{proof}
\begin{definition}[Matches Primary Pressure Volume Figure]
  \label{def:physics:phyx-mini-0492:phyxminiproblems-problemphyxmini0492-matchesprimarypressurevolumefigure}
  \lean{PhyXMiniProblems.ProblemPhyXMini0492.MatchesPrimaryPressureVolumeFigure}
  \uses{def:physics:phyx-mini-0492:phyxminiproblems-problemphyxmini0492-pressureinatmospheres, def:physics:phyx-mini-0492:phyxminiproblems-problemphyxmini0492-volumeinliters, def:physics:phyx-mini-0492:phyxminiproblems-problemphyxmini0492-twostepidealgasprocess}
  Exact transcription of the supplied pressure--volume bitmap. The red curve passes through a and c; the actual process instead follows the vertical arrow a → b and the horizontal arrow b → c
\end{definition}
\begin{proof}
  This is the declared model definition of Matches Primary Pressure Volume Figure.
\end{proof}
\begin{definition}[Has Physical Thermodynamic Parameters]
  \label{def:physics:phyx-mini-0492:phyxminiproblems-problemphyxmini0492-hasphysicalthermodynamicparameters}
  \lean{PhyXMiniProblems.ProblemPhyXMini0492.HasPhysicalThermodynamicParameters}
  \uses{def:physics:phyx-mini-0492:phyxminiproblems-problemphyxmini0492-pressureinatmospheres, def:physics:phyx-mini-0492:phyxminiproblems-problemphyxmini0492-volumeinliters, def:physics:phyx-mini-0492:phyxminiproblems-problemphyxmini0492-twostepidealgasprocess}
  Positivity and nondegeneracy conditions for the physical process
\end{definition}
\begin{proof}
  This is the declared model definition of Has Physical Thermodynamic Parameters.
\end{proof}
\begin{definition}[Satisfies Boundary Work Laws]
  \label{def:physics:phyx-mini-0492:phyxminiproblems-problemphyxmini0492-satisfiesboundaryworklaws}
  \lean{PhyXMiniProblems.ProblemPhyXMini0492.SatisfiesBoundaryWorkLaws}
  \uses{def:physics:phyx-mini-0492:phyxminiproblems-problemphyxmini0492-pressureinatmospheres, def:physics:phyx-mini-0492:phyxminiproblems-problemphyxmini0492-volumeinliters, def:physics:phyx-mini-0492:phyxminiproblems-problemphyxmini0492-energyinjoules, def:physics:phyx-mini-0492:phyxminiproblems-problemphyxmini0492-joulesperliteratmosphere, def:physics:phyx-mini-0492:phyxminiproblems-problemphyxmini0492-twostepidealgasprocess}
  General boundary-work laws for the two process kinds, written in the named units of the diagram. A constant-volume leg does no boundary work, while a constant-pressure expansion does P ΔV work. These laws contain none of the problem's numerical coordinates or answer values
\end{definition}
\begin{proof}
  This is the declared model definition of Satisfies Boundary Work Laws.
\end{proof}
\begin{definition}[Answer Choice]
  \label{def:physics:phyx-mini-0492:phyxminiproblems-problemphyxmini0492-answerchoice}
  \lean{PhyXMiniProblems.ProblemPhyXMini0492.AnswerChoice}
  Labels printed beside the four multiple-choice answers
\end{definition}
\begin{proof}
  This is the declared model definition of Answer Choice.
\end{proof}
\begin{definition}[work In Joules]
  \label{def:physics:phyx-mini-0492:phyxminiproblems-problemphyxmini0492-answerchoice-workinjoules}
  \lean{PhyXMiniProblems.ProblemPhyXMini0492.AnswerChoice.workInJoules}
  \uses{def:physics:phyx-mini-0492:phyxminiproblems-problemphyxmini0492-answerchoice}
  Work in joules printed beside each answer label
\end{definition}
\begin{proof}
  This is the declared model definition of work In Joules.
\end{proof}
\begin{definition}[recorded Answer Choice]
  \label{def:physics:phyx-mini-0492:phyxminiproblems-problemphyxmini0492-recordedanswerchoice}
  \lean{PhyXMiniProblems.ProblemPhyXMini0492.recordedAnswerChoice}
  \uses{def:physics:phyx-mini-0492:phyxminiproblems-problemphyxmini0492-answerchoice}
  Answer label recorded by the supplied dataset metadata
\end{definition}
\begin{proof}
  This is the declared model definition of recorded Answer Choice.
\end{proof}
\begin{definition}[Rounds To Nearest Ten Joules]
  \label{def:physics:phyx-mini-0492:phyxminiproblems-problemphyxmini0492-roundstonearesttenjoules}
  \lean{PhyXMiniProblems.ProblemPhyXMini0492.RoundsToNearestTenJoules}
  The answer choices report work to the nearest ten joules
\end{definition}
\begin{proof}
  This is the declared model definition of Rounds To Nearest Ten Joules.
\end{proof}
% --- Archon declaration topology end ---
% --- Archon physics formalization source end ---
